% ====================================================================
% CNT_LINK_DRAFT.tex  —  v1.0.22 R2 phase / 창 O-C (통계역학 증축) / 후보 (iv) 경량안
%   초안 전용 — 마스터 소스 미수정. 삽입처 = _sections/ch1_sec04_hys.tex
%   §4.3 (sec:hys-branch, 방향별 분기 중심), γ_j 지위 문단 직후(현행 251행 뒤).
%   목적 = §4 과주행(준안정 가지) 서술 ↔ 상분리 부록의 고전 핵생성 이론(CNT)의
%          최소 연결. ★신규 유도 금지 — 부록에 이미 완비(축의 연결 문단만).
%
%   ── 참조 방식 결정(브리핑 지시대로 마스터 tex 실확인 후) ──────────────────
%   마스터 ch1_graphite_v1.0.22.tex 는 appendix_phase_separation.tex 를 \input 하지
%   않는다(부록 4본만 input: ch1_appA/appB · ch2_appA/appB). 상분리 부록은 독립
%   컴파일 문서(자체 .pdf/.aux/.log)이며 \externaldocument 대상도 아니다.
%   ⇒ 참조는 \ref/xr 이 아니라 ★서술형(부록 명칭·절 제목 텍스트)로 유지한다.
%      (브리핑: "포함 아니면 서술형 참조 유지 — xr 대상 아님".)
%
%   ── 부록 라벨 실확인(현행 appendix_phase_separation.tex) ──────────────────
%   · 핵생성(CNT: ΔG*=16πγ³/(3Δg_v²)·핵생성률 ∝exp(-ΔG*/k_BT))
%       = 절 "준안정 영역 --- 핵생성"  label=app:nucleation  → 현행 §A.6.1
%         (식 eq:app-cnt · eq:app-rstar)
%   · 본문 연결 서술 = 절 "본문과의 연결"  label=app:link  → 현행 §A.7
%   ★주의: 브리핑/DIRECTION 이 적은 "§A.7(CNT)·§A.8/§A.9(연결)" 는 v1.0.20 구조 기준의
%      옛 번호로, 현행 부록(§A.6.1 핵생성·§A.7 연결)과 어긋난다. 서술형 참조라 절 제목으로
%      적어 번호 변동에 견고하게 둔다(REMOVAL_CHECK.md 에 기록).
%
%   ── 가드(DIRECTION 후보 (iv) 위험) ────────────────────────────────────────
%   · γ_j 현상학 지위 불변 — CNT 는 γ_j 의 "근거"이지 예측식이 아니다(모델 차원 불변).
%   · 기호 배향: 부록 ξ = Li 점유(=본문 θ), 본문 ξ = 진행률(1-θ) — 여집합.
%   · 기호 충돌: 부록 γ = 계면 에너지[J/m²] ≠ 여기 분기 축소 인자 γ_j (§7 이미 명시 구분).
%     → 아래 초안은 ΔG*·exp(-ΔG*/k_BT) 만 부르고 γ³/ξ-의존식은 부르지 않아 두 함정을 회피.
%   설계 원전 = v1.0.20 DIRECTION_STATMECH_REPORT.md 후보 (iv) 경량안.
%   사용 서지 = 신규 0(부록 [A3] Porter–Easterling·[A5] Balluffi 재사용 — 부록 소관).
% ====================================================================


% ─────────────────────────────────────────────────────────────────────
% 연결 문단  (§4.3 γ_j 지위 문단 "...넘을 수 없는 열역학 상한을 준다." 직후 삽입)
%   2~5문장 최소 연결. \ref 없음(서술형) · 신규 라벨 없음 · 신규 식 없음.
% ─────────────────────────────────────────────────────────────────────

이 조기 이탈의 미시 기작은 상분리 부록(독립 문서 --- \emph{상분리의 열역학})의 고전 핵생성
이론(classical nucleation theory, CNT)이 정량한다(부록 ``준안정 영역 --- 핵생성'' 절). 준안정
가지 위에서 새 상의 임계핵 형성에는 계면 비용이 세운 자유에너지 장벽 $\Delta G^\ast$ 이 있어
핵생성률이 $\propto\exp(-\Delta G^\ast/\kB T)$ 로 지수 억제되고, binodal(공통접선) 접점 근방에서는
구동력이 $0$ 으로 가 $\Delta G^\ast$ 가 발산하여 분해가 사실상 멈춘다 --- 그 지수 억제가 걷히기
전까지 소인이 전위를 계속 옮겨 가지를 spinodal 쪽으로 과주행시킨다(부록 ``본문과의 연결'' 절이 이
그림을 \S\ref{sec:hys} 과 잇는 반대 방향 서술을 이미 둔다). 곧 결함$\cdot$요동$\cdot$불균일이 없는
이상 단일 입자의 상한 $\gamma_j=1$(spinodal)에 대해, 실재 전극이 국소 요동$\cdot$불균일 핵생성으로
그 전에 가지를 이탈해 $\gamma_j<1$ 로 앉는 까닭에 CNT 가 미시 근거를 준다. 다만 CNT 는 $\gamma_j$ 의
\emph{근거}이지 예측식이 아니다 --- $\gamma_j$ 는 전이당 한 계수로 남는 현상학 인자이며(모델 차원
불변, spinodal gap $\Delta U_j^\hys$ 가 그 상한), 부록의 계면 에너지 $\gamma$[J/m$^2$]는 여기 축소
인자 $\gamma_j$ 와 무관한 별개 기호이고 부록의 자리 분율 $\xi$(=Li 점유=본문 $\theta$)도 본문 진행률
$\xi(=1-\theta)$ 와 여집합임에 유의한다.
